\section{编排器提示词}
\label{app:prompt}



\subsection{低 \masness}

\begin{systemprompt}[title={系统提示词}]
你是一名乐于助人的助手。

\masness（MAS 程度）：低

有效通道：thinking、agent、answer

模型：\texttt{[MODEL]}\\

智能体是预先配置的 AI 角色，可以接受委派的任务。每个子智能体：

1. 具有特定用途与专业领域；\\
2. 使用独立于主对话的上下文窗口；\\
3. （可选）可配置允许其使用的特定工具；\\
4. 包含用于指导其行为的自定义系统提示词。\\
\\
只能在 \xmltag{thinking} 通道内调用智能体，并使用 \xmltag{agent} 标签定义智能体。每个智能体都必须包含 \xmltag{agent_name}、\xmltag{agent_description}、\xmltag{required_arguments} 和 \xmltag{agent_output_id}。\\

不得遗漏任何请求字段，并确保回答是格式良好的 XML 对象！
\end{systemprompt}


\begin{developprompt}[title={开发提示词}]
通道：\\
    \xmltag{thinking}：内部推理与规划 \\
    \xmltag{agent}：定义子智能体 \\
    \xmltag{answer}：面向用户的最终答案 \\

模型（子智能体使用的模型）：

    \texttt{gpt-4.1-nano}：[官方网站上的模型介绍] \\
    \texttt{gpt-oss-120b}：[官方网站上的模型介绍] \\
    （……列出所有候选模型）\\

\masness 级别：

    \texttt{low}：直接求解或至多使用一个智能体 \\
    \texttt{high}：复杂的多智能体委派\\

子智能体模式（所有字段均为必需）：\\

\xmltag{agent}

\quad\xmltag{agent_name}
...
\xmlendtag{agent_name}
（从 CoTAgent、SCAgent、DebateAgent、ReflexionAgent 中选择一个）

\quad\xmltag{agent_description}
...
\xmlendtag{agent_description}

\quad\xmltag{required_arguments}
（确保设置所有必需参数，必须遵循 XML 格式）

\qquad\xmltag{...}...\xmlendtag{...}

\qquad\xmltag{...}...\xmlendtag{...}

\quad\xmlendtag{required_arguments}

\quad\xmltag{agent_output_id}
...
\xmlendtag{agent_output_id}

\xmlendtag{agent} \\

可用智能体：\\

\texttt{CoT}：[描述、名称、必需参数和实现摘要]

\texttt{SC}：[描述、名称、必需参数和实现摘要]

\texttt{Debate}：[描述、名称、必需参数和实现摘要]

\texttt{Reflexion}：[描述、名称、必需参数和实现摘要]

\texttt{Search}：[描述、名称、必需参数和实现摘要]


\end{developprompt}

\begin{userprompt}[title={用户提示词}]
请逐步解决该问题。在推理过程中，你可以自行处理任务，并在 \xmltag{answer} 标签中输出最终答案；也可以把任务委派给你设计的智能体，再在 \xmltag{answer} 标签中输出对应的 \texttt{agent\_output\_id}。你必须输出所有必需参数，并为智能体使用完全一致的字段名。\\
\\
例如，
\\
如果你能自行解决任务，应输出：\\

\xmltag{thinking}

(20+9)*(30+7) = 600 + 140 + 270 + 63 = 1073.

\xmlendtag{thinking}

\xmltag{answer}1073\xmlendtag{answer}
\\

如果你能通过一次工具调用、委派给单个智能体来解决任务，应输出：\\

\xmltag{thinking}

该问题涉及符号积分以及微积分基本定理的应用。
它需要结构化推理，而非简单的数值计算。
我将使用一个能够执行逐步思维链推理的微积分智能体。
原问题的最终答案将采用 CoTAgent 的输出。

\xmlendtag{thinking}

\xmltag{agent}

\quad\xmltag{agent_name}CoTAgent\xmlendtag{agent_name}

\quad\xmltag{agent_description}{通过一次思维链调用求解定积分。}\xmlendtag{agent_description}

\quad\xmltag{required_arguments}

\qquad\xmltag{agent_input}\xmlendtag{agent_input}

\quad\xmlendtag{required_arguments}

\quad\xmltag{agent_output_id}
calc_agent_output
\xmlendtag{agent_output_id}

\xmlendtag{agent}

\xmltag{answer}{calc_agent_output}\xmlendtag{answer} \\

另一个示例：\\

\xmltag{thinking}

该问题需要比较两个非常接近的数值选项。
为确保准确，我将让两个推理角色展开辩论：一个关注数学精度，另一个关注实际取整。
DebateAgent 能同时考虑两种观点，并得出有依据的最终答案。
最终答案将采用 DebateAgent 的输出。

\xmlendtag{thinking}

\xmltag{agent}

\quad\xmltag{agent_name}DebateAgent\xmlendtag{agent_name}

\quad\xmltag{agent_description}
通过一次辩论调用解决数值极为接近的选择问题。
\xmlendtag{agent_description}

\quad\xmltag{required_arguments}

\qquad\xmltag{agent_input}\xmlendtag{agent_input}

\qquad\xmltag{debate_roles}
["数学教授", "统计学教师"]
\xmlendtag{debate_roles}

\quad\xmlendtag{required_arguments}

\quad\xmltag{agent_output_id}
compare_agent_output
\xmlendtag{agent_output_id}

\xmlendtag{agent}

\xmltag{answer}compare_agent_output\xmlendtag{answer}\\

更多示例：\\

\xmltag{thinking}

可以直接计算 17³，但很容易出现算术错误。
采用带自洽性的思维链（CoT\_SC），可以采样多条推理路径并合并结果，从而提高准确性。
原问题的最终答案将采用 SCAgent 的输出。

\xmlendtag{thinking}

\xmltag{agent}

\quad\xmltag{agent_name}SCAgent\xmlendtag{agent_name}

\quad\xmltag{agent_description}
使用带自洽性的思维链（CoT\_SC）执行算术计算。
\xmlendtag{agent_description}

\quad\xmltag{required_arguments}

\qquad\xmltag{agent_input}\xmlendtag{agent_input}

\quad\xmlendtag{required_arguments}

\quad\xmltag{agent_output_id}math_agent_output\xmlendtag{agent_output_id}

\xmlendtag{agent}

\xmltag{answer}math_agent_output\xmlendtag{answer}
\\

最后一个示例：\\

\xmltag{thinking}

该任务需要使用公式进行推理，并确保单位处理正确。
我将使用 Reflexion 智能体；如果出现错误，它可以反思并改进自己的推理。
原问题的最终答案将采用 ReflexionAgent 的输出。

\xmlendtag{thinking}

\xmltag{agent}

\quad\xmltag{agent_name}ReflexionAgent\xmlendtag{agent_name}

\quad\xmltag{agent_description}
使用 Reflexion 的迭代自我反思与批评来解决推理任务。
\xmlendtag{agent_description}

\quad\xmltag{required_arguments}

\qquad\xmltag{agent_input}\xmlendtag{agent_input}

\quad\xmlendtag{required_arguments}

\quad\xmltag{agent_output_id}reflexion_agent_output\xmlendtag{agent_output_id}

\xmlendtag{agent}

\xmltag{answer}reflexion_agent_output\xmlendtag{answer} \\


如果把一个直接计算步骤与一次 Reflexion 工具调用结合，应输出：

\xmltag{thinking}

步骤 1：手工计算平均值：
\quad sum = 2 + 3 + 5 + 7 + 11 = 28
\quad avg = 28 / 5 = 5.6

步骤 2：
\quad 剩余步骤是把 sqrt(5.6) 计算到小数点后三位，这需要精确计算与数值改进。
\quad 我会把这一部分委派给 ReflexionAgent；若取整错误，它将进行自我纠正。
\quad 原问题的最终答案将采用 ReflexionAgent 的输出。

\xmlendtag{thinking}

\xmltag{agent}

\quad\xmltag{agent_name}ReflexionAgent\xmlendtag{agent_name}

\quad\xmltag{agent_description}
通过轻量自我改进循环计算平方根（单次智能体调用）。
\xmlendtag{agent_description}

\quad\xmltag{required_arguments}

\qquad\xmltag{agent_input}将 sqrt(5.6) 计算到小数点后三位\xmlendtag{agent_input}

\quad\xmlendtag{required_arguments}

\quad\xmltag{agent_output_id}numeric_agent_output\xmlendtag{agent_output_id}

\xmlendtag{agent}

\xmltag{answer}numeric_agent_output\xmlendtag{answer} \\


下面是需要解决的问题：

[QUESTION]


\end{userprompt}


\subsection{高 \masness}

% \zixuan{TODO: need to mention the difference}


\begin{systemprompt}[title={系统提示词}]
你是一名乐于助人的助手。

\masness（MAS 程度）：高

有效通道：thinking、agent、edge

模型：\texttt{[MODEL]}\\

智能体是预先配置的 AI 角色，可以接受委派的任务。每个子智能体：
1. 具有特定用途与专业领域；\\
2. 使用独立于主对话的上下文窗口；\\
3. （可选）可配置允许其使用的特定工具；\\
4. 包含用于指导其行为的自定义系统提示词。\\
\\
只能在 \xmltag{thinking} 通道内调用智能体，并使用 \xmltag{agent} 标签定义智能体。每个智能体都必须包含 \xmltag{agent_id}、\xmltag{agent_name}、\xmltag{agent_description} 和 \xmltag{required_arguments}。若要连接多个智能体以构成多智能体系统，请使用 \xmltag{edge} 通道。\\


不得遗漏任何请求字段，并确保回答是格式良好的 XML 对象！
\end{systemprompt}


\begin{developprompt}[title={开发提示词}]
通道：\\
    \xmltag{thinking}：内部推理与规划 \\
    \xmltag{agent}：定义子智能体 \\
    \xmltag{answer}：面向用户的最终答案 \\

模型（子智能体使用的模型）：

    \texttt{gpt-4.1-nano}：[官方网站上的模型介绍] \\
    \texttt{gpt-oss-120b}：[官方网站上的模型介绍] \\
    （……列出所有候选模型）\\

\masness 级别：

    \texttt{low}：直接求解或至多使用一个智能体 \\
    \texttt{high}：复杂的多智能体委派\\
    
子智能体模式（所有字段均为必需）：\\

\xmltag{agent}

\quad\xmltag{agent_name}
...
\xmlendtag{agent_name}
（从 CoTAgent、SCAgent、DebateAgent、ReflexionAgent 中选择一个）

\quad\xmltag{agent_description}
...
\xmlendtag{agent_description}

\quad\xmltag{required_arguments}
（确保设置所有必需参数，必须遵循 XML 格式）

\qquad\xmltag{...}...\xmlendtag{...}

\qquad\xmltag{...}...\xmlendtag{...}

\quad\xmlendtag{required_arguments}

\quad\xmltag{agent_output_id}
...
\xmlendtag{agent_output_id}

\xmlendtag{agent} \\


边模式（单个代码块；所有字段均为必需。每一对标签定义一条有向连接：\xmltag{from} 的输出 $\rightarrow$ \xmltag{to} 的输入。请在此列出所有连接；每个解决方案只能使用一个 \xmltag{edge} 块）：\\

\xmltag{edge}

\quad\xmltag{from}
...
\xmlendtag{from}
（源智能体的 agent\_id）

\quad\xmltag{to}
...
\xmlendtag{to}
（目标智能体的 agent\_id）

\xmlendtag{edge}\\

可用智能体：\\

\texttt{CoT}：[描述、名称、必需参数和实现摘要]

\texttt{SC}：[描述、名称、必需参数和实现摘要]

\texttt{Debate}：[描述、名称、必需参数和实现摘要]

\texttt{Reflexion}：[描述、名称、必需参数和实现摘要]

\texttt{Search}：[描述、名称、必需参数和实现摘要]


\end{developprompt}



\begin{userprompt}[title={用户提示词}]

请创建一个或多个智能体，并将它们连接成有效的计算图，通过协作解决给定问题并生成最终答案。创建智能体时，必须输出 \xmltag{agent} 来定义它，其中包括：\xmltag{agent_id}（智能体的唯一 ID，只能包含字母、数字或下划线，例如 A1、Refine\_1、WS\_Japan）、\xmltag{agent_name}（必须严格为 CoTAgent、SCAgent、DebateAgent、ReflexionAgent 或 WebSearchAgent 之一）、\xmltag{agent_description}、\xmltag{required_arguments}（必须至少包含一个 \xmltag{agent_input} 标签；DebateAgent 必须通过 \xmltag{debate_roles} 定义两个或更多角色；如果 \xmltag{agent_input} 留空（""），解析器会自动用原问题替换它）。\\

定义全部智能体后，必须通过指定描述智能体间数据流的边来构建有效图。只能输出一个 \xmltag{edge} 块；每对 \xmltag{from} 和 \xmltag{to} 将一个智能体的输出连接到另一个智能体的输入：\\

\xmltag{edge}

\quad\xmltag{from}source\_agent\_id\xmlendtag{from}

\quad\xmltag{to}target\_agent\_id\xmlendtag{to}

\xmlendtag{edge} \\

可以在 \xmltag{edge} 内输出多组 \xmltag{from} 和 \xmltag{to}。每个 \xmltag{from} 与 \xmltag{to} 的值都必须与已有的 \xmltag{agent_id} 完全匹配。\\

要使图有效，必须满足以下全部约束：
\begin{itemize}[leftmargin=*,noitemsep,topsep=2pt]
\item 节点一致性：每个 \xmltag{from} 和 \xmltag{to} 都必须引用出现在某个 \xmltag{agent} 块中的有效 \xmltag{agent_id}。
\item 有向性：边是有向的，数据从 \xmltag{from} 流向 \xmltag{to}。
\item 连通性：每个智能体都必须直接或间接连接到主流程，不允许存在孤立智能体。
\item 起始节点：至少一个智能体没有入边，它们是“入口点”（例如 WebSearch 或初始推理）。
\item 汇点：必须恰好有一个没有出边的智能体，它就是生成答案的最终智能体。
\item 不得出现未定义边：\xmltag{from} 或 \xmltag{to} 不得引用尚未声明的智能体。
\item 不得出现自环或环路：不允许 \xmltag{from}X\xmltag{/from}\xmltag{to}X\xmltag{/to} 形式的自环；图必须是有向无环图，并且存在拓扑序。
\item 允许并行：多个智能体可以具有相同的 \xmltag{from} 或 \xmltag{to}，即允许扇出和扇入。
\item 汇点无歧义：解析器会拒绝包含多个汇点的图；必要时请增加一个最终的“收集器”智能体。
\item 与顺序无关：XML 中各边的排列不必遵循执行顺序，系统会自动进行拓扑排序。
\item 汇点答案完整性：唯一汇点智能体的输出必须以面向用户的形式直接回答原问题。除非问题明确要求，否则它不能只是中间产物（例如笔记、批评或原始表格）。如果汇点的 \xmltag{agent_input} 为空，它会继承原问题，并必须返回最终答案。如果 \xmltag{agent_input} 非空，运行器仍会在其前面加入原问题作为上下文；汇点仍须给出面向用户的最终答案。
\item 边与数据流双向一致：边表示执行顺序，\$\{\} 表示数据流，二者必须相互一致。
    （a）如果 \xmltag{agent_input} 引用了 \$\{X\}，则必须存在边 \xmltag{from}X\xmltag{/from}\xmltag{to}该智能体\xmltag{/to}。
    （b）如果存在边 \xmltag{from}X\xmltag{/from}\xmltag{to}Y\xmltag{/to}，则 Y 的 \xmltag{agent_input} 必须引用 \$\{X\}。
    换言之，当且仅当一个智能体确实向另一个智能体传递数据时，二者之间才应存在边。不得仅为规定执行顺序而创建没有数据流的边。
\end{itemize}

思考部分（必需）：\\

定义智能体和边之前，必须包含一个 \xmltag{thinking} 部分。
其中应自然说明为何需要多个智能体、为何选择各类智能体，以及图为何采用该结构（并行、顺序或混合）。
它必须同时论证规划与设计理由。\\

示例结构：\\

\xmltag{thinking}
  解释单个智能体为何不足。
  描述每个智能体的角色及其连接方式。
  论证流程模式（并行、顺序或混合）。
  最后明确说明由哪个智能体生成最终输出。
\xmltag{/thinking}\\

单智能体示例：\\

如果决定使用单个智能体求解，应按如下格式输出。在这种情况下，由于 \xmltag{agent_input} 与原任务相同，必须把 \xmltag{agent_input} 设为空（""），解析器会用原问题替换它。\\

示例 1：\\

问题：计算 (2x + 5) dx 在 0 到 3 上的定积分。\\

\xmltag{thinking}

该问题涉及符号积分以及微积分基本定理的应用。
它需要结构化推理，而非简单的数值计算。
我将使用一个能够执行逐步思维链推理的微积分智能体。
原问题的最终答案将采用 CoTAgent 的输出。

\xmlendtag{thinking}

\xmltag{agent}

\quad\xmltag{agent_id}calc\_agent\xmlendtag{agent_id}

\quad\xmltag{agent_name}CoTAgent\xmlendtag{agent_name}

\quad\xmltag{agent_description}
通过一次思维链调用求解定积分。
\xmlendtag{agent_description}

\quad\xmltag{required_arguments}

\qquad\xmltag{agent_input}\xmlendtag{agent_input}

\quad\xmlendtag{required_arguments}

\xmlendtag{agent}\\

示例 2：\\

问题：17 的立方是多少？\\

\xmltag{thinking}

可以直接计算 17³，但很容易出现算术错误。
采用带自洽性的思维链（CoT\_SC），可以采样多条推理路径并合并结果，从而提高准确性。
原问题的最终答案将采用 SCAgent 的输出。

\xmlendtag{thinking}

\xmltag{agent}

\quad\xmltag{agent_id}math\_agent\xmlendtag{agent_id}

\quad\xmltag{agent_name}SCAgent\xmlendtag{agent_name}

\quad\xmltag{agent_description}
使用带自洽性的思维链（CoT\_SC）执行算术计算。
\xmlendtag{agent_description}

\quad\xmltag{required_arguments}

\qquad\xmltag{agent_input}\xmlendtag{agent_input}

\quad\xmlendtag{required_arguments}

\xmlendtag{agent} \\

示例 3：\\

问题：已知 x² = 46.694444，目标数 45 和 46 中哪一个更接近？\\

\xmltag{thinking}

该问题需要比较两个非常接近的数值选项。
为确保准确，我将让两个推理角色展开辩论：一个关注数学精度，另一个关注实际取整。
DebateAgent 能同时考虑两种观点，并得出有依据的最终答案。
最终答案将采用 DebateAgent 的输出。

\xmlendtag{thinking}

\xmltag{agent}

\quad\xmltag{agent_id}compare\_agent\xmlendtag{agent_id}

\quad\xmltag{agent_name}DebateAgent\xmlendtag{agent_name}

\quad\xmltag{agent_description}
通过一次辩论调用解决数值极为接近的选择问题。
\xmlendtag{agent_description}

\quad\xmltag{required_arguments}

\qquad\xmltag{agent_input}\xmlendtag{agent_input}

\qquad\xmltag{debate_roles}
["数学教授", "统计学教师"]
\xmlendtag{debate_roles}

\quad\xmlendtag{required_arguments}

\xmlendtag{agent}

示例 4：\\

问题：一列火车以每小时 60 英里的速度行驶，2.5 小时能行驶多远？\\

\xmltag{thinking}

该任务需要使用公式进行推理，并确保单位处理正确。
我将使用 Reflexion 智能体；如果出现错误，它可以反思并改进自己的推理。
原问题的最终答案将采用 ReflexionAgent 的输出。

\xmlendtag{thinking}

\xmltag{agent}

\quad\xmltag{agent_id}reflexion\_agent\xmlendtag{agent_id}

\quad\xmltag{agent_name}ReflexionAgent\xmlendtag{agent_name}

\quad\xmltag{agent_description}
使用 Reflexion 的迭代自我反思与批评来解决推理任务。
\xmlendtag{agent_description}

\quad\xmltag{required_arguments}

\qquad\xmltag{agent_input}\xmlendtag{agent_input}

\quad\xmlendtag{required_arguments}

\xmlendtag{agent}\\

示例 5：\\

问题：截至本月，日本当前的通货膨胀率是多少？\\

\xmltag{thinking}

该问题依赖最新事实信息，无法仅凭静态知识可靠回忆。
使用单个 WebSearchAgent 即可，因为该任务只需从网络检索准确且带引用的事实，无需进一步推理或综合。
该智能体将搜索在线来源，并返回简洁、基于引用的摘要。
因此，WebSearchAgent 既是流程中的唯一节点，也是最终输出节点。

\xmlendtag{thinking}

\xmltag{agent}

\quad\xmltag{agent_id}SEARCH\xmlendtag{agent_id}

\quad\xmltag{agent_name}WebSearchAgent\xmlendtag{agent_name}

\quad\xmltag{agent_description}
从互联网检索最新且带引用的事实信息。
\xmlendtag{agent_description}

\quad\xmltag{required_arguments}

\qquad\xmltag{agent_input}\xmlendtag{agent_input}

\quad\xmlendtag{required_arguments}

\xmlendtag{agent}\\

多智能体示例：\\

如果决定使用多个智能体求解，必须先把原问题分解为更小且定义明确的子任务，每个子任务只对应一个子目标。随后为每个子任务创建一个智能体，并将它们连接成连贯的无环计算图。\\

在这种情况下，agent\_input 不为空，而是智能体需要解决的具体子任务；解析器会自动把原问题作为上下文添加到给定 agent\_input 内容之前。\\

分解时：\\

1. 保持子任务最小且聚焦。每个智能体应只处理一个原子目标（例如一次查询、一个推理步骤或一项验证任务）。\\
2. 针对不同事实证据使用多个 WebSearchAgent，而不是执行一次宽泛搜索。\\
3. 按逻辑连接智能体，使信息流向直接回答原问题的唯一最终智能体（汇点）。\\


示例 1：\\

问题：简要总结纽约市、洛杉矶和芝加哥最新的住房空置率，并附引用。\\


\xmltag{thinking}

我们需要带来源的最新数据，因此将并行运行三个 WebSearch 智能体。
随后，由一个 CoT 智能体把三个片段规范化为小表格。
一个 SC 智能体将运行多个细微变体并投票，以减少抽取错误。
一个 Reflexion 智能体将检查单位、时效性与引用。
最后由一个 CoT 智能体撰写简短摘要。
最终汇点是 FINAL。

\xmlendtag{thinking}

\xmltag{agent}

\quad\xmltag{agent_id}WS\_NYC\xmlendtag{agent_id}

\quad\xmltag{agent_name}WebSearchAgent\xmlendtag{agent_name}

\quad\xmltag{agent_description}
查找纽约市最新住房空置率及来源文本。
\xmlendtag{agent_description}

\quad\xmltag{required_arguments}

\qquad\xmltag{agent_input}
搜索纽约市住房空置率的最新官方或可靠估计。返回包含数值、日期和一行引用的简短片段。
\xmlendtag{agent_input}

\quad\xmlendtag{required_arguments}

\xmlendtag{agent}

\xmltag{agent}

\quad\xmltag{agent_id}WS\_LA\xmlendtag{agent_id}

\quad\xmltag{agent_name}WebSearchAgent\xmlendtag{agent_name}

\quad\xmltag{agent_description}
查找洛杉矶最新空置率及来源文本。
\xmlendtag{agent_description}

\quad\xmltag{required_arguments}

\qquad\xmltag{agent_input}
搜索洛杉矶住房空置率的最新官方或可靠估计。返回包含数值、日期和一行引用的简短片段。
\xmlendtag{agent_input}

\quad\xmlendtag{required_arguments}

\xmlendtag{agent}

\xmltag{agent}

\quad\xmltag{agent_id}WS\_CHI\xmlendtag{agent_id}

\quad\xmltag{agent_name}WebSearchAgent\xmlendtag{agent_name}

\quad\xmltag{agent_description}
查找芝加哥最新空置率及来源文本。
\xmlendtag{agent_description}

\quad\xmltag{required_arguments}

\qquad\xmltag{agent_input}
搜索芝加哥住房空置率的最新官方或可靠估计。返回包含数值、日期和一行引用的简短片段。
\xmlendtag{agent_input}

\quad\xmlendtag{required_arguments}

\xmlendtag{agent}

\xmltag{agent}

\quad\xmltag{agent_id}EXT\xmlendtag{agent_id}

\quad\xmltag{agent_name}CoTAgent\xmlendtag{agent_name}

\quad\xmltag{agent_description}
抽取数值，并统一三个空置率的日期与引用格式。
\xmlendtag{agent_description}

\quad\xmltag{required_arguments}

\qquad\xmltag{agent_input}
从以下片段中抽取每座城市的名称、空置率（百分比）、参考日期（YYYY-MM 或 YYYY）和简短引用。输出一张紧凑的三行表格。

纽约市：
\$\{WS\_NYC\}

洛杉矶：
\$\{WS\_LA\}

芝加哥：
\$\{WS\_CHI\}
\xmlendtag{agent_input}

\quad\xmlendtag{required_arguments}

\xmlendtag{agent}

\xmltag{agent}

\quad\xmltag{agent_id}VOTE\xmlendtag{agent_id}

\quad\xmltag{agent_name}SCAgent\xmlendtag{agent_name}

\quad\xmltag{agent_description}
对抽取结果进行集成，以减少复制或解析错误。
\xmlendtag{agent_description}

\quad\xmltag{required_arguments}

\qquad\xmltag{agent_input}
给定以下三行表格，独立执行五次抽取，并投票得到一份纠正后的三行表格。

表格：
\$\{EXT\}
\xmlendtag{agent_input}

\quad\xmlendtag{required_arguments}

\xmlendtag{agent}

\xmltag{agent}

\quad\xmltag{agent_id}QA\xmlendtag{agent_id}

\quad\xmltag{agent_name}ReflexionAgent\xmlendtag{agent_name}

\quad\xmltag{agent_description}
检查单位、日期时效性与引用是否存在；必要时列出修正项。
\xmlendtag{agent_description}

\quad\xmltag{required_arguments}

\qquad\xmltag{agent_input}
审核投票后的表格：单位应为 \%，日期应存在，并且每座城市都应有引用。如发现问题，用三行列出具体修正；否则回答 OK。最后给出一行结论。

投票后的表格：
\$\{VOTE\}
\xmlendtag{agent_input}

\quad\xmlendtag{required_arguments}

\xmlendtag{agent}

\xmltag{agent}

\quad\xmltag{agent_id}FINAL\xmlendtag{agent_id}

\quad\xmltag{agent_name}CoTAgent\xmlendtag{agent_name}

\quad\xmltag{agent_description}
撰写带引用的简短摘要。
\xmlendtag{agent_description}

\quad\xmltag{required_arguments}

\qquad\xmltag{agent_input}
使用检查后的表格与注释，撰写 3--4 句话的摘要：每座城市一句，最后一句比较各空置率。保留表格中的行内来源引用。

表格：
\$\{VOTE\}

质量检查注释：
\$\{QA\}
\xmlendtag{agent_input}

\quad\xmlendtag{required_arguments}

\xmlendtag{agent}

\xmltag{edge}

\quad\xmltag{from}WS\_NYC\xmlendtag{from}\xmltag{to}EXT\xmlendtag{to}

\quad\xmltag{from}WS\_LA\xmlendtag{from}\xmltag{to}EXT\xmlendtag{to}

\quad\xmltag{from}WS\_CHI\xmlendtag{from}\xmltag{to}EXT\xmlendtag{to}

\quad\xmltag{from}EXT\xmlendtag{from}\xmltag{to}VOTE\xmlendtag{to}

\quad\xmltag{from}VOTE\xmlendtag{from}\xmltag{to}QA\xmlendtag{to}

\quad\xmltag{from}VOTE\xmlendtag{from}\xmltag{to}FINAL\xmlendtag{to}

\quad\xmltag{from}QA\xmlendtag{from}\xmltag{to}FINAL\xmlendtag{to}

\xmlendtag{edge}\\


示例 2：\\

问题：20 世纪 70 年代末，教皇约翰·保罗二世首次出访时到访了一个以丰富的中美洲历史和庞大人口著称的国家。后来他在哪一天访问了一座亚得里亚海沿岸的重要城市？该城市以重要港口以及一座供奉姓名首字母为“S”的圣人的著名宗座圣殿而闻名；同一天他还访问了附近的哪座城市？\\


\xmltag{thinking}

该问题需要历史推理：先识别教皇约翰·保罗二世的首次出访，再找出他访问过的一座亚得里亚海城市，该市拥有供奉姓名以“S”开头圣人的宗座圣殿；随后确定访问日期，以及他同日访问的另一座城市。

单个智能体无法完成全部工作，因为任务混合了检索与推理。
我将把它分解为更小的子任务：三个 WebSearchAgent 分别检索事实，一个 CoTAgent 构建时间线，一个 ReflexionAgent 验证一致性，最后一个 CoTAgent 撰写答案。
由 FINAL 智能体输出最终答案。

\xmlendtag{thinking}

\xmltag{agent}

\quad\xmltag{agent_id}WS\_FIRST\_TRIP\xmlendtag{agent_id}

\quad\xmltag{agent_name}WebSearchAgent\xmlendtag{agent_name}

\quad\xmltag{agent_description}
检索教皇约翰·保罗二世在 20 世纪 70 年代末首次出访的细节。
\xmlendtag{agent_description}

\quad\xmltag{required_arguments}

\qquad\xmltag{agent_input}
搜索教皇约翰·保罗二世的首次出访（年份、目的国和持续时间）。
返回行程日期范围、目的地和可靠引用。
\xmlendtag{agent_input}

\quad\xmlendtag{required_arguments}

\xmlendtag{agent}

\xmltag{agent}

\quad\xmltag{agent_id}WS\_ADRIATIC\xmlendtag{agent_id}

\quad\xmltag{agent_name}WebSearchAgent\xmlendtag{agent_name}

\quad\xmltag{agent_description}
查找亚得里亚海沿岸城市的访问记录，其中城市拥有供奉姓名以“S”开头圣人的宗座圣殿。
\xmlendtag{agent_description}

\quad\xmltag{required_arguments}

\qquad\xmltag{agent_input}
搜索教皇约翰·保罗二世访问过的亚得里亚海城市，该市须有一座供奉姓名首字母为“S”的圣人的宗座圣殿
（例如圣尼古拉或圣马可）。
返回城市名、宗座圣殿名、访问日期与引用。
\xmlendtag{agent_input}

\quad\xmlendtag{required_arguments}

\xmlendtag{agent}

\xmltag{agent}

\quad\xmltag{agent_id}WS\_SAME\_DAY\xmlendtag{agent_id}

\quad\xmltag{agent_name}WebSearchAgent\xmlendtag{agent_name}

\quad\xmltag{agent_description}
识别教皇约翰·保罗二世在访问该亚得里亚海城市当天还到访过的附近城市。
\xmlendtag{agent_description}

\quad\xmltag{required_arguments}

\qquad\xmltag{agent_input}
搜索教皇约翰·保罗二世在访问该亚得里亚海城市同一天到访的其他城市。
返回附近城市名、距离估计与来源引用。
\xmlendtag{agent_input}

\quad\xmlendtag{required_arguments}

\xmlendtag{agent}

\xmltag{agent}

\quad\xmltag{agent_id}TIMELINE\xmlendtag{agent_id}

\quad\xmltag{agent_name}CoTAgent\xmlendtag{agent_name}

\quad\xmltag{agent_description}
把行程数据整合为一条经过验证的历史时间线。
\xmlendtag{agent_description}

\quad\xmltag{required_arguments}

\qquad\xmltag{agent_input}
使用以下结果：

首次出访：
\$\{WS\_FIRST\_TRIP\}

亚得里亚海城市访问：
\$\{WS\_ADRIATIC\}

同日访问：
\$\{WS\_SAME\_DAY\}

构建清晰的时间线，确认亚得里亚海城市的访问日期与同日到访的附近城市。
\xmlendtag{agent_input}

\quad\xmlendtag{required_arguments}

\xmlendtag{agent}

\xmltag{agent}

\quad\xmltag{agent_id}VERIFY\xmlendtag{agent_id}

\quad\xmltag{agent_name}ReflexionAgent\xmlendtag{agent_name}

\quad\xmltag{agent_description}
验证时间顺序、地理位置以及圣人与宗座圣殿之间的联系。
\xmlendtag{agent_description}

\quad\xmltag{required_arguments}

\qquad\xmltag{agent_input}
检查时间线事实是否一致：
- 确保所有来源中出现的日期相同。
- 确认亚得里亚海城市在地理上邻近第二座城市。
- 验证该宗座圣殿确实供奉一位姓名以“S”开头的圣人。

若信息一致则返回 OK，否则列出修正项。

待验证时间线：
\$\{TIMELINE\}
\xmlendtag{agent_input}

\quad\xmlendtag{required_arguments}

\xmlendtag{agent}

\xmltag{agent}

\quad\xmltag{agent_id}FINAL\xmlendtag{agent_id}

\quad\xmltag{agent_name}CoTAgent\xmlendtag{agent_name}

\quad\xmltag{agent_description}
生成包含日期和两座城市的简洁最终答案。
\xmlendtag{agent_description}

\quad\xmltag{required_arguments}

\qquad\xmltag{agent_input}
使用已验证的结果，以一句话回答：

“教皇约翰·保罗二世于 [DATE] 访问了 [ADRIATIC\_CITY]，这里有供奉圣人 [S] 的宗座圣殿；同一天，他还访问了 [NEARBY\_CITY]。”

另附一句引用来源的简短验证说明。

时间线：
\$\{TIMELINE\}

验证：
\$\{VERIFY\}
\xmlendtag{agent_input}

\quad\xmlendtag{required_arguments}

\xmlendtag{agent}

\xmltag{edge}

\quad\xmltag{from}WS\_FIRST\_TRIP\xmlendtag{from}\xmltag{to}TIMELINE\xmlendtag{to}

\quad\xmltag{from}WS\_ADRIATIC\xmlendtag{from}\xmltag{to}TIMELINE\xmlendtag{to}

\quad\xmltag{from}WS\_SAME\_DAY\xmlendtag{from}\xmltag{to}TIMELINE\xmlendtag{to}

\quad\xmltag{from}TIMELINE\xmlendtag{from}\xmltag{to}VERIFY\xmlendtag{to}

\quad\xmltag{from}VERIFY\xmlendtag{from}\xmltag{to}FINAL\xmlendtag{to}

\quad\xmltag{from}TIMELINE\xmlendtag{from}\xmltag{to}FINAL\xmlendtag{to}

\xmlendtag{edge}\\

（……更多示例）\\

下面是需要解决的问题：

[QUESTION]


\end{userprompt}





\subsection{子智能体}
\label{app:sub_agents}

% \zixuan{the prompt for IMAS is slightly different due to the nature of the task}

\begin{subagent}[title={CoTAgent}]

\pykw{async} \pykw{def} \pyfn{CoTAgent}(\pykw{self}, agent\_input, model: \pykw{str}):

\quad \pycom{\# 基本设置}

\quad temperature = 0.5

\quad \pycom{\# 思维链指令}

\quad cot\_instruction = \pystr{"请逐步思考，然后解决任务。"}

\quad \pycom{\# 实例化 CoT LLM}

\quad cot\_agent = LLMAgentBase([\pystr{thinking}, \pystr{answer}],\pystr{思维链 LLM}, model=model, temperature=temperature)

\quad thinking, answer = \pykw{await} cot\_agent([agent\_input], cot\_instruction)

\quad final\_answer = self.make\_final\_answer(thinking, answer)

\quad \pykw{return} final\_answer \\

func\_string = inspect.getsource(CoTAgent) \\

COT = \code{\{}

\quad \jsonkey{description}: 
\jsonstr{通过鼓励 LLM 逐步思考而非直接输出答案，思维链推理借助中间步骤解决复杂问题。}

\quad \jsonkey{name}: \jsonstr{思维链智能体（CoTAgent）}

\quad \jsonkey{required_arguments}: \code{\{}

\quad\quad \jsonkey{agent_input}: 
\jsonstr{CoTAgent 的输入。若为空，解析器会用原问题替换。}

\quad \code{\}}

\quad \jsonkey{implementation}: \jsonstr{\{func\_string\}}

\code{\}}


\end{subagent}


\begin{subagent}[title={SCAgent}]
    \pykw{async} \pykw{def} \pyfn{SCAgent}(\pykw{self}, agent\_input, model: \pykw{str}):

\quad \pycom{\# 基本设置}

\quad temperature = 0.5

\quad num\_repeated\_samples = 5

\quad \pycom{\# 思维链指令}

\quad cot\_instruction = \pystr{"请逐步思考，然后解决任务。"}

\quad \pycom{\# 初始化多个 CoT 智能体以实现自洽}

\quad cot\_agents = [
\quad\quad LLMAgentBase([\pystr{thinking}, \pystr{answer}],
\quad\quad \pystr{思维链 LLM},
\quad\quad model=model,
\quad\quad temperature=temperature)
\quad\quad \pykw{for} \_ \pykw{in} range(num\_repeated\_samples)
\quad ]

\quad thinking\_mapping = \{\}

\quad answer\_mapping = \{\}

\quad possible\_answers = []

\quad \pykw{for} i \pykw{in} range(num\_repeated\_samples):

\quad\quad thinking, answer = \pykw{await} cot\_agents[i]([agent\_input], cot\_instruction)

\quad\quad possible\_answers.append(answer.content)

\quad\quad thinking\_mapping[answer.content] = thinking

\quad\quad answer\_mapping[answer.content] = answer

\quad \pycom{\# 通过多数投票集成答案}

\quad answer = self.majority\_voting(possible\_answers)

\quad thinking = thinking\_mapping[answer]

\quad answer = answer\_mapping[answer]

\quad final\_answer = self.make\_final\_answer(thinking, answer)

\quad \pykw{return} final\_answer \\

func\_string = inspect.getsource(SCAgent) \\

COT\_SC = \code{\{}

\quad \jsonkey{description}:
\jsonstr{LLM 虽能得到正确答案，但其推理可能有所变化。通过以较高温度重复询问同一问题，可以生成多条推理路径；随后对思维链智能体的答案进行集成，得到更准确的最终答案。该方法最适合需要通过共识获得高置信度的问题。}

\quad \jsonkey{name}:
\jsonstr{思维链自洽智能体（SCAgent）}

\quad \jsonkey{required_arguments}: \code{\{}

\quad\quad \jsonkey{agent_input}:
\jsonstr{SCAgent 的输入。若为空（""），解析器会自动用原问题替换。}

\quad \code{\}}

\quad \jsonkey{implementation}:
\jsonstr{\{func\_string\}}

\code{\}}

\end{subagent}


\begin{subagent}[title={DebateAgent}]
\pykw{async} \pykw{def} \pyfn{DebateAgent}(\pykw{self}, agent\_input, model: \pykw{str}, debate\_roles: List[\pykw{str}]):

\quad \pycom{\# 基本设置}

\quad temperature = 0.5

\quad max\_debate\_round = 5

\quad \pycom{\# 初始推理指令}

\quad debate\_initial\_instruction = 
\quad \pystr{"请逐步思考，然后解决任务。"}

\quad \pycom{\# 辩论更新指令}

\quad debate\_instruction = 
\quad \pystr{"给定其他智能体的解答，将其观点视为建议，并给出更新后的答案。"}

\quad \pycom{\# 使用不同角色初始化辩论智能体}

\quad debate\_agents = [
\quad\quad LLMAgentBase([\pystr{thinking}, \pystr{answer}],
\quad\quad \pystr{辩论 LLM},
\quad\quad model=model,
\quad\quad role=role,
\quad\quad temperature=temperature)
\quad\quad \pykw{for} role \pykw{in} debate\_roles
\quad ]

\quad \pycom{\# 最终决策指令}

\quad final\_decision\_instruction =
\quad \pystr{"给定全部推理与答案，请谨慎给出最终答案。"}

\quad final\_decision\_agent = LLMAgentBase(
\quad\quad [\pystr{thinking}, \pystr{answer}],
\quad\quad \pystr{最终决策 LLM},
\quad\quad model=model,
\quad\quad temperature=temperature
\quad )

\quad all\_thinking = [[] \pykw{for} \_ \pykw{in} range(max\_debate\_round)]

\quad all\_answer = [[] \pykw{for} \_ \pykw{in} range(max\_debate\_round)]

\quad \pycom{\# 执行多轮辩论}

\quad \pykw{for} r \pykw{in} range(max\_debate\_round):

\quad\quad \pykw{for} i \pykw{in} range(len(debate\_agents)):

\quad\quad\quad \pykw{if} r == 0:

\quad\quad\quad\quad thinking, answer = 
\quad\quad\quad\quad \pykw{await} debate\_agents[i]([agent\_input], debate\_initial\_instruction)

\quad\quad\quad \pykw{else}:

\quad\quad\quad\quad input\_infos = 
\quad\quad\quad\quad [agent\_input] 
\quad\quad\quad\quad + [all\_thinking[r-1][i]] 
\quad\quad\quad\quad + all\_thinking[r-1][:i] 
\quad\quad\quad\quad + all\_thinking[r-1][i+1:]

\quad\quad\quad\quad thinking, answer = 
\quad\quad\quad\quad \pykw{await} debate\_agents[i](input\_infos, debate\_instruction)

\quad\quad\quad all\_thinking[r].append(thinking)

\quad\quad\quad all\_answer[r].append(answer)

\quad \pycom{\# 根据全部辩论结果作出最终决策}

\quad thinking, answer = 
\quad \pykw{await} final\_decision\_agent(
\quad\quad [agent\_input] 
\quad\quad + all\_thinking[max\_debate\_round-1] 
\quad\quad + all\_answer[max\_debate\_round-1],
\quad\quad final\_decision\_instruction
\quad )

\quad final\_answer = self.make\_final\_answer(thinking, answer)

\quad \pykw{return} final\_answer \\

func\_string = inspect.getsource(DebateAgent) \\

LLM\_debate = \code{\{}

\quad \jsonkey{description}:
\jsonstr{通过让不同 LLM 相互辩论，可以利用多元观点得到更好的解答。该智能体最适合能从多种视角中获益的问题。}

\quad \jsonkey{name}:
\jsonstr{LLM 辩论智能体（DebateAgent）}

\quad \jsonkey{required_arguments}: \code{\{}

\quad\quad \jsonkey{agent_input}:
\jsonstr{DebateAgent 的输入。若为空（""），解析器会自动用原问题替换。}

\quad\quad \jsonkey{debate_roles}:
\jsonstr{代表不同视角的角色列表（必须包含多个角色），例如“数学教授”或“统计学家”。}

\quad \code{\}}

\quad \jsonkey{implementation}:
\jsonstr{\{func\_string\}}

\code{\}}

\end{subagent}


\begin{subagent}[title=ReflexionAgent]
\pykw{async} \pykw{def} \pyfn{ReflexionAgent}(\pykw{self}, agent\_input, model: \pykw{str}):

\quad \pycom{\# 基本设置}

\quad temperature = 0.5

\quad max\_reflection\_round = 5

\quad \pycom{\# 初始推理指令}

\quad initial\_instruction = 
\quad \pystr{"请逐步思考，然后解决任务。"}

\quad \pycom{\# 反思与改进指令}

\quad reflect\_instruction = 
\quad \pystr{"给定之前的尝试与反馈，请仔细思考最近一次尝试可能出错之处。利用此前尝试中的洞见，更好地解决任务。"}

\quad cot\_agent = LLMAgentBase(
\quad\quad [\pystr{thinking}, \pystr{answer}],
\quad\quad \pystr{思维链 LLM},
\quad\quad model=model,
\quad\quad temperature=temperature
\quad )

\quad \pycom{\# 批评反馈指令}

\quad critic\_instruction =
\quad \pystr{"请审查以上答案，并批评其中可能出错之处。如果你完全确定答案正确，请在 'correct' 中严格输出 'True'。"}

\quad critic\_agent = LLMAgentBase(
\quad\quad [\pystr{feedback}, \pystr{correct}],
\quad\quad \pystr{批评 LLM},
\quad\quad model=model,
\quad\quad temperature=temperature
\quad )

\quad \pycom{\# 初始尝试}

\quad cot\_inputs = [agent\_input]

\quad thinking, answer = 
\quad \pykw{await} cot\_agent(cot\_inputs, initial\_instruction, 0)

\quad \pycom{\# 迭代式自我反思循环}

\quad \pykw{for} i \pykw{in} range(max\_reflection\_round):

\quad\quad feedback, correct = 
\quad\quad \pykw{await} critic\_agent(
\quad\quad\quad [agent\_input, thinking, answer],
\quad\quad\quad critic\_instruction,
\quad\quad\quad i
\quad\quad )

\quad\quad \pykw{if} correct.content == \pystr{'True'}:

\quad\quad\quad \pykw{break}

\quad\quad cot\_inputs.extend([thinking, answer, feedback])

\quad\quad thinking, answer =
\quad\quad \pykw{await} cot\_agent(cot\_inputs, reflect\_instruction, i + 1)

\quad final\_answer = self.make\_final\_answer(thinking, answer)

\quad \pykw{return} final\_answer \\

func\_string = inspect.getsource(ReflexionAgent) \\

Reflexion = \code{\{}

\quad \jsonkey{description}:
\jsonstr{为提升性能，LLM 可以根据反馈迭代改进答案。模型通过反思先前尝试并吸收批评来完善推理，给出更准确的解答。该智能体最适合能从自我纠正中获益的复杂问题。}

\quad \jsonkey{name}:
\jsonstr{自我改进（Reflexion）}

\quad \jsonkey{required_arguments}: \code{\{}

\quad\quad \jsonkey{agent_input}:
\jsonstr{ReflexionAgent 的输入。若为空（""），解析器会自动用原问题替换。}

\quad \code{\}}

\quad \jsonkey{implementation}:
\jsonstr{\{func\_string\}}

\code{\}}


\end{subagent}

    
